% 原书第 160—161 页；第 10.5 节。
\section{讨论}
\label{sec:ch10-discussion}

在 van der Laan and Gruber（2010）中，我们提出了构造迭代算法的一般模板。该算法建立一个有序的 TMLE 序列
$(Q_{n,j}^*,G_{n,j}^*)$，$j=1,\ldots,J$，使相关部分 $Q_{n,j}$ 和干扰参数 $G_{n,j}$ 的经验拟合随 $j$ 增大而改善，再使用基于损失的交叉验证，选择 $Q_0$ 的最佳估计量 $Q_{n,j_n}^*$。本章假设已经给定 $G_0$ 的这样一个有序估计量集合，这对应于上述一般 C-TMLE 模板的一个特例：一般模板通过评价一组可能的改动，在 $G_{n,j}$ 的基础上构造下一个估计量 $G_{n,j+1}$；在本章的特例中，下一个 $G_{n,j+1}$ 已知，即只有一种改动。

不过，我们研究了指标参数 $h$ 连续取值的情形，从而超越以往关于 C-TMLE 的文献。我们关注由二阶余项 $R_{20}(Q,G,Q_0,G_0)$ 的乘积结构所定义的双重稳健估计问题，并假设候选估计量族 $\{G_{n,h}:h\}$ 允许一致估计 $G_0$。我们遵循一般 C-TMLE 模板，描述了一个略作修改的 C-TMLE 算法，以保证所选估计量 $G_{n,h_n}$ 求解调节参数 $h$ 的得分方程。我们说明，求解这个得分方程，意味着 TMLE $(Q_{n,h_n}^*,G_{n,h_n})$ 求解一个至关重要的得分方程
\[
 P_n\frac{d}{dh_n}D^*(Q_n^*,G_{n,h_n})\approx0,
\]
其中，这个得分定义为典范梯度对 $h$ 的导数。我们解释了为什么后一个得分方程，对应于局部极小化 TMLE 的 Taylor 展开中的二阶余项
$R_{20}(Q_n^*,G_{n,h},Q_0,G_0)$，这里的极小化是对 $h$ 而言。

此外，我们证明了一个正式定理：除了通常的有效影响曲线方程
$P_nD^*(Q_n^*,G_{n,h_n})=0$
之外，再求解这一关键方程，就能保证在可能的情况下，选择器 $h_n$ 使估计足够欠平滑，从而即使 $Q_0$ 的初始估计量不一致或收敛缓慢，TMLE $\Psi(Q_n^*)$ 仍然保持渐近线性。另一方面，采用交叉验证选择器 $h_{n,\mathrm{CV}}$、仅求解
$P_nD^*(Q_n^*,G_{n,h_{n,\mathrm{CV}}})=0$
的 TMLE，在 $Q_n^*$ 不一致或收敛缓慢时，一般不具有渐近线性。我们还说明，使用一个以二维 $\epsilon$ 生成这两个得分的最不利子模型，可以得到同时求解两个得分方程的“特殊”TMLE。于是，即使采用 $h_{n,\mathrm{CV}}$，相同的渐近理论也适用于这一标准 TMLE。不过，基于 C-TMLE 在有限样本中的理论依据，我们预期 C-TMLE 的有限样本表现会优于这种特殊 TMLE。这里没有展示的初步模拟表明，C-TMLE 和特殊 TMLE 都很容易优于标准 TMLE；在“标准 TMLE”“特殊 TMLE”和“C-TMLE”这三个估计量中，C-TMLE 的表现似乎最好。

